% !TeX root = ../main.tex
\begin{abstract}
Small language models can be trained to conduct Motivational Interviewing (MI) by
simulating sessions against an LLM patient and rewarding them with an LLM that fills in
validated MI questionnaires. We ask what such training actually buys. Holding the base
model (Llama-3.2-1B), the patient simulator, the grading oracle, the candidate budget and
the look-ahead depth fixed, we compare two ways of turning the same reward into a policy
update over ten iterations: Preference-Tree Optimization (PTO), which grows a sequential
search over therapist turns and trains with DPO on the surviving pairs, and Group-Relative
Policy Optimization (GRPO), which trains on standardized advantages over a group of
completions. Every model state is evaluated on 96 patient personas across eight
instruments.
%
Both methods improve substantially on the rubric they optimize, but PTO leads at the
matched ten-iteration endpoint (\QQ{} $4.26$ vs $3.75$; paired $+0.51$, $\dz{}=0.73$),
because GRPO peaks at iteration~8 and then regresses. The gains arrive with a measurable
drift toward affirmation: literal question rates fall from $0.83$ to $0.15$ per therapist
turn in GRPO, and MI-inconsistent behaviour rises ${\sim}2.3\times$ (PTO) and
${\sim}4\times$ (GRPO).
%
Re-scoring all $22{,}272$ conversation$\times$rubric cells with a held-out judge from a
different model family preserves $88.3\%$ of contrast signs ($98.9\%$ at
$|\Delta|\!\ge\!0.50$), but separates the arms on \emph{gain retention}: the fraction of
each arm's Q1 gain the held-out judge still credits diverges from iteration~3 onward,
ending at $0.80$ for PTO and $0.28$ for GRPO. Finally, reading the training signal
directly shows the difference is a property of the \emph{data} the two methods collect,
not of their losses. We argue that questionnaire-filling LLM judges should be treated as
optimization targets that must themselves be audited, and that gain retention under a
decoupled grader is a cheap, general test for doing so.
\end{abstract}
